%% degrees-and-priority-continuation.tex
%% Independently revised continuation of historical body topics 152--435.

% Necessity: the exact inherited source ends before the later degree,
% finite-injury, and c.e.-structure material already used by downstream parts.
% Definition: this is an independently revised continuation of the distinct
% degree and priority topics preserved in historical body lines 152--435.
% Standing: the continuation is editorial handbook content, not inherited
% restoration or a claim that the historical derivative remains byte exact.
% Nonredundancy: the exact source supplies oracle computation, jump, and Post's
% theorem, while this continuation begins with Turing degrees and omits that core.
% Use: the Part IV driver routes it once after all exact source through line
% 9079, and later handbook consumers retain its degree and priority results.
\noindent\textit{Reader note.}
The following separately authored continuation retains the nonoverlapping
degree and priority material from the preserved historical draft.  It follows
the exact inherited relative-computability chapter and does not replace or
alter that source presentation.

\section{Turing degrees}
\label{sec:p04-turing-degrees}

\begin{hbdefinition}{def:p04-turing-degree}{Turing equivalence and degrees}
The preorder \(\leq_T\) compares sets but does not identify different sets
with the same computational power.  Define
\[
 A\sim_TB\quad\Leftrightarrow\quad
 A\leq_TB\ \text{ and }\ B\leq_TA.
\]
The \emph{Turing degree} of \(A\) is its equivalence class
\(\deg_T(A)=[A]_T\).  The set of degrees is ordered by
\[
 [A]_T\leq[B]_T\quad\Leftrightarrow\quad A\leq_TB.
\]
This is well defined and antisymmetric precisely because representatives are
quotiented by mutual reducibility.  A degree therefore names an equivalence
class of sets, not a set, a machine, or a reducibility statement.  The
distinction is used by join, jump, and Post's problem below.
\index{Turing degree}
\index{Turing equivalence}
\end{hbdefinition}

\begin{hblemma}{thm:p04-jump-extensive}{proved=annex:p4.degrees.1}{Every oracle reduces to its jump}
For every \(A\subseteq\N\),
\[
 A\leq_m A'.
\]
\end{hblemma}

Together with the nonreducibility half of the jump theorem, this gives
\(\deg_T(A)<\deg_T(A)'\).  The lower reduction is proved in the annex; the
diagonal nonreducibility theorem and its complete source are
\cite{KleenePost1954,Rogers1967}.

\begin{hbdefinition}{def:p04-degree-operations}{Join and jump on degrees}
For \(A,B\subseteq\N\), put
\[
 A\oplus B=\{2n:n\in A\}\cup\{2n+1:n\in B\}.
\]
Define the join of \(\deg_T(A)\) and \(\deg_T(B)\) to be
\(\deg_T(A\oplus B)\).  Lemma~\ref{thm:p04-degree-semilattice} verifies that
this operation is independent of representatives and is their least upper
bound.  The monotonicity in
Remark~\ref{rem:jump-convention} makes
\[
 \deg_T(A)'\defeq\deg_T(A')
\]
well defined.  Write \(\mathbf0=\deg_T(\varnothing)\) and
\(\mathbf0'=\deg_T(\varnothing')\).  A \emph{c.e. degree} is the degree of
some computably enumerable set.  These are derived operations and subclasses
of the already defined degree order, not additional primitive relations.
\index{Turing degree!join}
\index{Turing degree!jump}
\index{computably enumerable degree}
\end{hbdefinition}

\begin{hblemma}{thm:p04-degree-semilattice}{cited=KleenePost1954}{Upper-semilattice structure of the Turing degrees}
The order on Turing degrees and the join operation in
Definition~\ref{def:p04-degree-operations} are independent of representatives.
The degree \(\deg_T(A\oplus B)\) is an upper bound of
\(\deg_T(A)\) and \(\deg_T(B)\), and it is below every common upper bound.
\end{hblemma}

\begin{hbmethod}{def:p04-kleene-post-construction}{Finite-extension construction below the jump}
Use the fixed oracle-program numbering \((\varphi_e^B)_{e<\omega}\) and meet
 \[
 R_e:\chi_A\not\simeq\varphi_e^B,
 \qquad
 S_e:\chi_B\not\simeq\varphi_e^A.
\]
Construct increasing finite characteristic strings
\(\sigma_s,\tau_s\in2^{<\omega}\), whose unions define \(A\) and \(B\),
using \(\varnothing'\) to decide the c.e. question
whether some finite extension of the opposite oracle makes the selected
functional converge with output \(0\) on a fresh coordinate.  If such an
extension exists, adopt it and put \(1\) into the target string; otherwise put
\(0\) there.  Alternate the \(R_e\)- and \(S_e\)-stages and extend both
strings in length at every stage.

This method is needed because the theorem requires actual incomparable sets
below \(\mathbf0'\), not merely a cardinality argument.  It is a
\(\varnothing'\)-effective construction of two characteristic functions, not
the later c.e. enumeration with restraints.  The negative oracle question is
what distinguishes it from an ordinary computable finite extension.  The
annex verifies totality, the two diagonal requirements, reducibility to
\(\varnothing'\), and strictness below \(\mathbf0'\).
\index{Kleene--Post construction}
\end{hbmethod}

\begin{hbtheorem}{thm:p04-kleene-post}{proved=annex:p4.degrees.2}{Kleene--Post incomparability theorem}
There are Turing degrees \(\mathbf a,\mathbf b<\mathbf0'\) such that neither
\(\mathbf a\leq\mathbf b\) nor \(\mathbf b\leq\mathbf a\).  No claim that
these two degrees are c.e. is made by this theorem.
\end{hbtheorem}

\section{Post's problem and finite injury}
\label{sec:p04-finite-injury}

Post's problem asks whether some c.e. degree lies strictly between
\(\mathbf0\) and \(\mathbf0'\).  Friedberg and Muchnik answered it
affirmatively by constructing incomparable c.e. sets
\cite{Friedberg1957,Muchnik1956}.

\begin{hbmethod}{def:p04-finite-injury-method}{Finite-injury priority construction}
Construct c.e. sets \(A\) and \(B\) in effective stages while meeting, in the
order
\[
 R_0>S_0>R_1>S_1>\cdots,
\]
the requirements
\[
 R_e:\varphi_e^B\not\simeq\chi_A,
 \qquad
 S_e:\varphi_e^A\not\simeq\chi_B.
\]
Here \(\varphi_e^B\) is the partial function computed by the \(e\)-th
program in the fixed acceptable oracle-machine numbering, while
\(\chi_A\) and \(\chi_B\) are the total
characteristic functions of the constructed sets.  The requirements therefore
compare functions extensionally rather than comparing a function with a set.
An \(R_e\)-strategy reserves a fresh witness \(x\) above all active
higher-priority restraints, above the current stage, and distinct from every reserved witness, keeps
\(x\notin A\), and waits for a convergent computation
\(\varphi_e^B(x)=0\).  Its query use \(u\) is the largest oracle position
queried by that finite computation; if it makes no query, its oracle restraint
is empty.  When the computation appears, enumerate \(x\) into \(A\), preserve
every oracle answer below \(u+1\), and declare the requirement satisfied.
Only a value \(0\) triggers action.  If the candidate functional were total
and equal to the target characteristic function while the final witness stayed
out, its stable value there would have to be \(0\), so the wait would
eventually end.  The
\(S_e\)-strategy is symmetric.

If a higher-priority strategy later enumerates below a lower strategy's
oracle restraint, initialize the lower strategy again: cancel its witness and
restraint, but never remove a number already enumerated; then choose a new
fresh witness.  No lower-priority strategy may enumerate within a protected
oracle segment unless a higher-priority action first injures and initializes
it.  This is a construction scheme for enumerations of \(A,B\), not a new
reducibility or a new kind of degree.  Annex~\ref{annex:p4.degrees.4} gives
the effective scheduler, finite-injury induction, and satisfaction proof.
\index{finite-injury priority method}
\index{priority argument}
\end{hbmethod}

\begin{hbtheorem}{thm:p04-friedberg-muchnik}{proved=annex:p4.degrees.4}{Friedberg--Muchnik theorem}
There are computably enumerable sets \(A,B\subseteq\N\) such that
\(A\not\leq_TB\) and \(B\not\leq_TA\).  Consequently there are incomparable
c.e. degrees strictly between \(\mathbf0\) and \(\mathbf0'\), answering
Post's problem.
\end{hbtheorem}

\begin{table}[htbp]
\centering
\small
\begin{tabular}{@{}p{.65in}p{1.15in}p{1.25in}p{1.25in}p{1.30in}@{}}
\toprule
Strategy & Witness policy & Wait condition & Positive action & Restraint and injury \\
\midrule
\(R_e\) & reserve fresh \(x\notin A\) & \(\varphi_e^B(x)\downarrow=0\) & enumerate \(x\) into \(A\) & preserve every \(B\)-answer below \(u+1\); reset after higher-priority injury \\
\(S_e\) & reserve fresh \(y\notin B\) & \(\varphi_e^A(y)\downarrow=0\) & enumerate \(y\) into \(B\) & preserve every \(A\)-answer below \(u+1\); reset after higher-priority injury \\
no action & retain the final witness outside the target & no stable value \(0\) appears & none & equality would force a stable value \(0\), so the requirement is already met \\
verification & induct on priority & higher strategies have stabilized & final witness is chosen & only finitely many resets; final restraint survives \\
\bottomrule
\end{tabular}
\caption{Finite-injury construction table for Friedberg--Muchnik.  The table
records the construction, while
Theorem~\ref{thm:p04-friedberg-muchnik} and its annex item carry the complete
proof.}
\label{tab:p04-finite-injury}
\end{table}

\section{Structure of the c.e. sets}
\label{sec:p04-ce-structure}

\begin{hbconvention}{def:p04-ce-numbering}{Acceptable numbering of the c.e. sets}
The index-sensitive results below require more than an arbitrary effective
enumeration.  For this chapter, an \emph{acceptable numbering of programs}
means a numbering satisfying the enumeration/surjectivity clause~(b) of
Theorem~\ref{thm:enumeration}, the universal property of
Theorem~\ref{thm:universal-computation}, and the uniform parameter
property of Theorem~\ref{thm:parameter}.  The fixed program numbering has
exactly this package.  Put
\[
 W_e\defeq\{x\in\N:\varphi_e(x)\mathord\downarrow\}.
\]
Thus \((W_e)_{e\in\N}\) is the induced acceptable numbering of the c.e.
subsets of \(\N\).  The adjective \emph{acceptable} here names a
metalinguistic package of numbering hypotheses, not a reducibility, a degree,
or a new class of sets; unlike a merely effective enumeration, it supplies
the uniform index translations used in productivity, creativity, Myhill's
theorem, and Rice--Shapiro.
\index{acceptable numbering}
\index{computably enumerable set!acceptable numbering}
\end{hbconvention}

\begin{hbdefinition}{def:p04-one-one-reduction}{One-one reducibility}
A \emph{one-one reduction} of \(A\subseteq\N\) to \(B\subseteq\N\) is a
total computable injective function \(f\) such that
\[
 x\in A\quad\Leftrightarrow\quad f(x)\in B.
\]
Write \(A\leq_1B\).  This strengthens many-one reducibility by requiring the
translator to be injective; the distinction is used by Myhill's theorem.
\index{one-one reducibility}
\end{hbdefinition}

\begin{hbdefinition}{def:p04-simple-set}{Simple set}
A c.e. set \(S\) is \emph{simple} when \(\N\setminus S\) is infinite but
contains no infinite c.e. subset.  Simplicity is a property of a c.e. set,
not a completeness notion: its use is to witness a noncomputable c.e. set.
Simplicity alone neither implies nor excludes Turing completeness.
\index{simple set}
\end{hbdefinition}

\begin{hbtheorem}{thm:p04-simple-set-exists}{cited=Post1944}{Existence of simple sets}
There exists a simple computably enumerable set.  Every simple set is
noncomputable.
\end{hbtheorem}

\begin{hbdefinition}{def:p04-productive-creative}{Productive complements and creative sets}
A set \(C\) is \emph{productive} when some partial computable function \(p\)
has the following property: whenever \(W_e\subseteq C\), the value \(p(e)\)
is defined and lies in \(C\setminus W_e\).  A c.e. set \(A\) is
\emph{creative} when its complement is productive.
Productivity supplies a uniform escape from every c.e. subset; creativity
combines that property with enumerability and is used in the completeness
theorem below.
\index{productive set}
\index{creative set}
\end{hbdefinition}

\begin{hbtheorem}{thm:p04-creative-complete}{cited=Myhill1955}{Creative-set characterization}
A computably enumerable set is creative if and only if it is
many-one complete among the computably enumerable sets.  In particular, the
self-application domain set \(K\) is creative.
\end{hbtheorem}

\begin{hbtheorem}{thm:p04-myhill}{cited=Myhill1955}{Myhill isomorphism theorem}
If \(A\leq_1B\) and \(B\leq_1A\), then some computable permutation
\(\pi:\N\to\N\) satisfies
\[
 x\in A\quad\Leftrightarrow\quad\pi(x)\in B.
\]
Consequently, any two creative sets are computably isomorphic.
\end{hbtheorem}

\begin{hbtheorem}{thm:p04-rice-shapiro}{cited=Johnson1977Rice}{Rice--Shapiro theorem}
For \(d\in\N\), let \(D_d\) be the finite subset of \(\N\) whose members are
the positions of the \(1\)-bits in the binary expansion of \(d\).  Let
\(\mathcal C\) be a class of c.e. subsets of \(\N\), and let
\(I_{\mathcal C}=\{e:W_e\in\mathcal C\}\) be its extensional index set.
Then \(I_{\mathcal C}\) is computably enumerable if and only if there is a
computably enumerable set \(\mathsf{Fin}_{\mathcal C}\subseteq\N\) of
finite-set codes such
that, for every c.e. set \(W\),
\[
 W\in\mathcal C
 \quad\Leftrightarrow\quad
 \exists d\in\mathsf{Fin}_{\mathcal C}\;(D_d\subseteq W).
\]
The empty class is included by allowing
\(\mathsf{Fin}_{\mathcal C}=\varnothing\).
Thus a positive c.e. extensional property is witnessed by finite positive
information.  The result concerns extensional index sets of c.e. sets, not
arbitrary semantic properties of programs; within that scope it is stronger
than Rice's undecidability classification.
\end{hbtheorem}

\begin{table}[htbp]
\centering
\small
\begin{tabular}{@{}p{1.12in}p{1.25in}p{1.30in}p{1.55in}@{}}
\toprule
Object/fact & Exact standing & Order-theoretic content & Evidence \\
\midrule
Turing degrees & \(\mathcal P(\N)/{\sim_T}\) & partial order induced by \(\leq_T\) & Definition~\ref{def:p04-turing-degree} \\
join & \(\deg_T(A\oplus B)\) & least upper bound of two degrees & Lemma~\ref{thm:p04-degree-semilattice}; \cite{KleenePost1954} \\
jump & \(\deg_T(A)'=\deg_T(A')\) & strictly increases every degree & Lemma~\ref{thm:p04-jump-extensive} and Theorem~\ref{thm:jump} \\
Kleene--Post & arbitrary degrees below \(\mathbf0'\) & incomparable pair exists & Theorem~\ref{thm:p04-kleene-post}; \cite{KleenePost1954} \\
Post's problem & c.e. degrees between \(\mathbf0\) and \(\mathbf0'\) & asks for a nontrivial intermediate c.e. degree & \cite{Post1944} \\
Friedberg--Muchnik & c.e. degrees & incomparable pair, hence intermediate degrees & complete annex proof; \cite{Friedberg1957,Muchnik1956} \\
\bottomrule
\end{tabular}
\caption{Turing degrees and jump facts.  The word
\emph{arbitrary} in the
Kleene--Post row prevents the false inference that its witnesses are the c.e.
witnesses later produced by finite injury.}
\label{tab:p04-degree-facts}
\end{table}

\begin{figure}[htbp]
\centering
\begin{tikzpicture}[
  node distance=14mm and 18mm,
  degree/.style={circle,draw,minimum size=9mm,font=\small},
  order/.style={-{Latex[length=2mm]},thick},
  jumpop/.style={-{Latex[length=2mm]},thick,dashed}
]
 \node[degree] (zero) {\(\mathbf0\)};
 \node[degree,above left=of zero] (a) {\(\mathbf a\)};
 \node[degree,above right=of zero] (b) {\(\mathbf b\)};
 \node[degree,above=28mm of zero] (jump) {\(\mathbf0'\)};
 \draw[order] (zero) -- (a);
 \draw[order] (zero) -- (b);
 \draw[order] (a) -- (jump);
 \draw[order] (b) -- (jump);
 \draw[jumpop,bend left=38] (zero) to
   node[right,font=\scriptsize] {jump \(\mathbf0\mapsto\mathbf0'\)} (jump);
 \node[font=\scriptsize] at ($(a)!0.5!(b)$) {incomparable};
\end{tikzpicture}
\caption{A schematic fragment of the degree order, not to scale.  Here
\(\mathbf a,\mathbf b\) are the c.e. degrees supplied
by Friedberg--Muchnik; the Kleene--Post theorem separately supplies an
unrestricted incomparable pair.  Solid arrows record only the displayed
strict comparisons; the dashed arrow is the established jump instance
\(\mathbf0\mapsto\mathbf0'\).  The figure asserts neither covers nor
distances, exhaustion of an interval, or the existence of a meet.}
\label{fig:p04-turing-degrees}
\end{figure}
